\clearpage
\SourceLeaf{274}%
\begin{center}
\textbf{[METAPHYSIK DER GEOMETRIE.]}

\medskip
[\textsc{W. Sartorius von Waltershausen.} \emph{Gauss zum Gedächtniss. Leipzig 1856. S. 80---81.}]
\end{center}

\{\mbox{.\,.\,.\,.\,.\,.\,.\,.\,.} In seiner frühesten Jugend habe ihm [\textsc{Gauss}] die Geometrie wenig Interesse eingeflößt, das sich bei ihm erst später in hohem Maß entwickelt habe.

Es war besonders bemerkenswert und überaus lehrreich zu sehen, wie \textsc{Gauss} die Grundlagen der Mathematik offenlegte und sie scharf von der Metaphysik abgrenzte. Obwohl er über diese Fragen nie etwas veröffentlicht hat, lässt sich vermuten, dass sich darüber in seinem wissenschaftlichen Nachlass einiges finden wird. In früherer Zeit, als sein Lebensweg noch offen war und er damit rechnen musste, vielleicht irgendwo als Lehrer der Mathematik zu arbeiten, hatte er für diese Möglichkeit eine Schrift ausgearbeitet, die noch in seinen letzten Jahren vorhanden gewesen sein soll und in der er die Anfänge der Mathematik philosophisch entwickelt hatte. Ob sie sich jetzt noch finden wird, ist zweifelhaft.

Die Geometrie betrachtete \textsc{Gauss} erst dann als ein folgerichtiges Gebäude, wenn man die Theorie der Parallelen als oberstes Axiom anerkenne. Er sei jedoch zu der Überzeugung gelangt, dass sich dieser Satz nicht beweisen lasse. Aus der Erfahrung, zum Beispiel aus den Winkeln des Dreiecks Brocken, Hohehagen, Inselsberg, wisse man aber, dass er näherungsweise richtig sei. Wollte man das genannte Axiom dagegen nicht anerkennen, so folge daraus eine andere, ganz\space%
